Setzen wir zweitens: $\tau=\sigma-\lambda$, $\rho=n-\sigma-\lambda$; $B^\sigma=A^\sigma$, so kommt:
\begin{equation}
\begin{aligned}
\pair{q_{n-\sigma-\lambda}A^\sigma}{A_{n-\sigma}p_{\sigma-\lambda}}
&=(-1)^\lambda\pair{q_{n-\sigma-\lambda}A^\sigma}{A_{n-\sigma}p_{\sigma-\lambda}}\\
&\quad+(-1)^{(n-\sigma)(\sigma-\lambda)}
 \sum_{\alpha=1}^{\sigma-\lambda,n-\sigma-\lambda}
 \Pi\pair{q_{n-\sigma-\lambda-\alpha}A^\sigma}{A_{n-\sigma}p_{\sigma-\lambda}}.
\end{aligned}
\tag{39}
\end{equation}
Es sind also, wie in \S\ 2 bemerkt, die die Symbolreihen charakterisierenden Bedingungen (12.) f\"ur ungeraden Defekt $\lambda$ eine Folge der Bedingungen von h\"oherem Defekt. F\"ur eine Linearform $\pair{A^\sigma}{p_\sigma}=\pair{B^\sigma}{p_\sigma}$ ergibt sich aus (39.), da\ss\ sich ihre irreduzibeln invarianten Bildungen, die in den Koeffizienten quadratisch sind, auf solche von geradem Defekt reduzieren; das stimmt mit der bekannten Tatsache \"uberein, da\ss\ eine Linearform im bin\"aren und tern\"aren Gebiet keine invarianten Bildungen besitzt.

Bemerkt sei, da\ss\ die allgemeinen Identit\"aten urspr\"unglich abgeleitet sind unter der Annahme, da\ss\ die einzelnen Reihen den Bedingungen (12.) gen\"ugen. Da indes diese einzelnen Reihen nur linear in die Identit\"aten eingehen, so gelten letztere auch f\"ur die allgemeineren, den Bedingungen (12.) nicht gen\"ugenden Reihen.
